\chapter{Conclusion}
\label{chap:conclusion}

This thesis set out to close a specific, narrow gap: the existing FAIR
Jupyter reproducibility pipeline can detect that a notebook fails because of
a missing or incompatible Python package, but it does not explain that
failure or attempt to repair it. \Cref{chap:related-work} showed that the
individual pieces needed to close this gap -- failure detection, automated
dependency inference, and LLM-based program repair -- already existed
separately in the literature, but not combined inside a reproducibility
pipeline, grounded in live package data, and paired with a plain-language
explanation.

\Cref{chap:problem-dataset} through \Cref{chap:validation} reported what has
actually been built to close it. Of the 214 dependency-error notebooks in
the Docker pipeline's own residual failures, 200 fall inside this thesis's
pip-only repair scope. A deterministic classifier decides scope and subtype
before any language model is involved. An explanation component turns a raw
error message into a plain-language, schema-validated description of what
went wrong, for every one of the 214 rows, not only the repairable ones. A
repair component proposes a fix using an open, locally-run model whose
suggestion is checked, field by field, against live PyPI metadata and a
small, hand-curated compatibility registry, so that a suggested package name
or version is never a guess. A fix-application component applies that
proposal inside a rebuilt environment equivalent to the one the notebook
originally failed in, re-runs the whole notebook, and classifies the result
into exactly three outcomes. All of this is implemented, covered by 310
automated tests, and additionally checked against real PyPI, a real Ollama
model, a real Docker daemon, and real GitHub repositories in a small pilot.

That pilot produced one finding worth restating here because it shapes how
the rest of this project should be read: every real case tried resolved its
targeted dependency error and then uncovered a second, previously invisible
one. This is not a weakness of the components built for this thesis --
their own classification logic is separately proven capable of reporting a
completely clean result -- but a real characteristic of research notebooks
that depend on more than one package. It is also a direct argument for the
two-round repair mode that is designed, but not yet built
(\Cref{chap:future-work}): a single repair attempt is often only the first
of several a real notebook may need.

Several design decisions changed along the way, and \Cref{chap:methodology}
recorded each one rather than hiding it behind a description that only
matches the final result. The clearest examples are that the original
containers this system was assumed to reuse for fix validation turned out
never to have been preserved, so the environment is rebuilt from the same
recipe instead; that the repair proposal's own persisted record does not
carry a notebook's identity forward, so a re-join step had to be added; and
that the explanation component's output format changed from free prose to
strict, schema-validated JSON once it had to fit inside an automated
pipeline. None of these changes weaken the system -- each was made because a
concrete implementation problem required it, and each is now covered by a
regression test.

What remains is stated plainly in \Cref{chap:future-work}: an orchestrator
that runs the existing components together over the full dataset without a
person driving each attempt by hand, persistence of results into the
designed \texttt{repair\_attempts} table and the FAIR Jupyter knowledge
graph, and the full-scale evaluation -- including a one-round versus
two-round comparison and a baseline comparison -- that would turn the pilot
observations in \Cref{chap:validation} into a measured repair success rate.
Reaching that point does not require redesigning any of the components
described in this thesis; it requires connecting components that are
already built, tested, and shown to work correctly on their own.
